\documentclass[../../main.tex]{subfiles}
\begin{document}

{\bf[解]}

\subparagraph{(1)}

% TODO: upgrade
\begin{figure}[htb]
  \centering
  \begin{tikzpicture}[scale=3, >=stealth, every node/.style={font=\small}]

    % --- 1. 座標定義 ---
    \coordinate (O)  at (0, 0);          % P_0 (原点)
    \coordinate (P1) at (0.5, 0);        % 1回目: 東へ 1/2
    \coordinate (P2) at (0.75, 0);       % 2回目: 東へ 1/4 (l-1回目まで x 軸上)
    \coordinate (P3) at (0.75, 0.40);    % l回目: 初めて北へ (1/2)^l (例として幅 0.40)
    \coordinate (P4) at (0.75, 0.20);    % l+1回目: 南へ (1/2)^{l+1} = 0.20
    \coordinate (P5) at (0.75, 0.10);    % l+2回目: 南へ (1/2)^{l+2} = 0.10
    \coordinate (Pn_lim) at (0.75, 0.05);% その後すべて南へ動いたときの極限位置 (> 0)

    % --- 2. 4つの象限の対称性ハイライト (2)用 ---
    \fill[blue!6, opacity=0.8] (0, 0) rectangle (1.4, 0.9);
    \node[blue!85!black, font=\footnotesize\bfseries] at (1.15, 0.72) {第1象限 ($a_n$)};
    \node[gray!60, font=\scriptsize] at (-0.6, 0.72) {第2象限};
    \node[gray!60, font=\scriptsize] at (-0.6, -0.6) {第3象限};
    \node[gray!60, font=\scriptsize] at (1.15, -0.6) {第4象限};

    % --- 3. 座標軸 ---
    \draw[->, thick, gray!80!black] (-0.9, 0) -- (1.5, 0) node[right] {$x$};
    \draw[->, thick, gray!80!black] (0, -0.8) -- (0, 1.0) node[above] {$y$};
    \node[below left=2pt] at (O) {$O(P_0)$};

    % --- 4. 点 P の移動パス ---
    % (i) 1回目〜l-1回目: x 軸上を移動 (東西)
    \draw[->, line width=1.5pt, teal!85!black] (O) -- (P1)
    node[midway, below=2pt, font=\scriptsize] {$\frac{1}{2}$};
    \draw[->, line width=1.5pt, teal!85!black] (P1) -- (P2)
    node[midway, below=2pt, font=\scriptsize] {$\dots$};

    % (ii) l 回目: 初めて南北（北）へ動く
    \draw[->, line width=2.0pt, red!85!black] (P2) -- (P3);

    % (iii) l+1 回目以降: 最大限 x 軸に向かって逆向き（南）に動くケース
    \draw[->, thick, dashed, orange!90!black] (P3) -- (P4);
    \draw[->, thick, dashed, orange!90!black] (P4) -- (P5);
    \draw[thick, dotted, orange!90!black] (P5) -- (Pn_lim);

    % --- 6. 各主要点のプロット ---
    \fill[black] (O) circle (1.0pt);
    \fill[teal!85!black] (P1) circle (1.2pt);
    \fill[teal!85!black] (P2) circle (1.2pt) node[below=2pt, font=\scriptsize] {$P_{l-1}$};
    \fill[red!85!black]  (P3) circle (1.5pt) node[above=2pt, font=\footnotesize] {$P_l$};
    \fill[orange!90!black] (Pn_lim) circle (1.3pt) node[above right=1pt, font=\scriptsize] {$P_n$};

  \end{tikzpicture}
  \caption{点 $P$ の移動軌跡の例}
  \label{fig:random_walk_axis_irreversibility}
\end{figure}

$l$回目$(a\le l\le n-1,l\in\mathbb{N})$にはじめて南北方向に動いたとすると，移動量は$\left(\dfrac{1}{2}\right)^l$である．ここで$l+1$回目から$n$回目まで全て$l$回目と逆向きに動いたとしても，その合計の移動量は
\begin{align}
  \left(\dfrac{1}{2}\right)^{l+1}+ \dots+\left(\dfrac{1}{2}\right)^n \\
  =\left(\dfrac{1}{2}\right)^l-\left(\dfrac{1}{2}\right)^n<\left(\dfrac{1}{2}\right)^l
\end{align}
だから，$P_n$は$x$軸上には存在しない．以上から題意の対偶が示された．
よって題意も示された．$\Box$

\subparagraph{(2)}

$P_n$が$x$軸上にあるという事象を$X$，$y$軸上にあるという事象を$Y$とする．
求める確率$a_n$とすると，対称性から他の象限にある確率も$a_n$であるから，
\begin{align}
  4a_n+P(X\cup Y)=1 \label{eq:1}
\end{align}
が成立する．従って以下$P(X\cup Y)$を計算することで$a_n$を求める．

前問の結果から$P_n$が軸上に有るとき$P_k(1\le k\le n)$も軸上にある．
$x$軸上にあるのは常にサイコロが$2,4,5,6$のとき，
$y$軸上にあるのは常にサイコロが$1,3,5,6$のときである．
したがって
\begin{align}
  P(X)=P(Y)=\left(\dfrac{2}{3}\right)^n  \label{eq:2}
\end{align}
である．また$P(X\cap Y)$は$P_n$が原点にある時で，
これは(1)と同様に考えれば常にサイコロが$5,6$であるときだから，
\begin{align}
  P(X\cap Y)=\left(\dfrac{1}{3}\right)^n \label{eq:3}
\end{align}
である．
\begin{align}
  P(X\cup Y)=P(X)+P(Y)-P(X\cap Y) \label{eq:4}
\end{align}
に\cref{eq:2,eq:3}を代入して
\begin{align}
  P(X\cup Y) = 2\left(\dfrac{2}{3}\right)^n - \left(\dfrac{1}{3}\right)^n
\end{align}
だから，これを\cref{eq:1}に代入して求める確率は
\begin{align}
  a_n
  =\dfrac{1}{4}\left(1-2\left(\dfrac{2}{3}\right)^n+\left(\dfrac{1}{3}\right)^n\right)\cdots\text{(答)}
\end{align}
となる．
\newpage
\end{document}
